% =============================================================================
%  List of Notations.
%
%  A short EXAMPLE list. Replace it with your own symbols: it is here to show
%  the shape and the ordering, not because these are the symbols you need.
%
%  Order entries the way a reader looks things up: operators first, then
%  letters, with the variants of one letter grouped together. Nothing generates
%  or checks this file, so a symbol renamed in a later chapter goes stale here
%  silently; keep it in step as you write.
%
%  Note \bm rather than \mathbf for bold maths, matching the notation rule in
%  \Cref{sec:notation}. The second argument to mclistof is the width reserved
%  for the terms.
% =============================================================================
\begin{mclistof}{List of Notations}{3.2cm}

    \item[$|\cdot|$] Measure of magnitude or size: absolute value for real numbers, norm for vectors, cardinality for sets.
    \item[$d$] Dimensionality of a vector.
    \item[$\eta$] Learning rate.
    \item[$\bm{h}$] Hidden state vector.
    \item[$\mathcal{L}$] Loss function.
    \item[$n$, $m$] Number of data points, samples or observations.
    \item[$\mathbb{N}$, $\mathbb{R}$] Sets of natural and of real numbers.
    \item[$p$] Probability.
    \item[$\bm{q}$, $\bm{k}$, $\bm{v}$] Query, key and value in the attention calculation \parencite{vaswani2017attention}.
    \item[$\bm{\theta}$] The parameters of a model, collectively.
    \item[$w$, $\bm{w}$, $\bm{W}$] Weights, for instance those of a neural network.
    \item[$x$, $\bm{x}$, $\bm{X}$] Input features, or independent variable.
    \item[$y$, $\bm{y}$, $\bm{Y}$] Target, or dependent variable.
    \item[$\hat{y}$, $\hat{\bm{y}}$] Predicted value of the target.

\end{mclistof}
